\documentclass[11pt]{article}

\usepackage[margin=1in]{geometry}
\usepackage[T1]{fontenc}
\usepackage[utf8]{inputenc}
\usepackage{lmodern}
\usepackage{amsmath,amssymb}
\usepackage{enumitem}
\usepackage{hyperref}
\usepackage{microtype}
\usepackage{graphicx}
\setlength{\parindent}{1.5em}
\setlength{\parskip}{0.35em}

\title{Information, Physics, Quantum: The Search for Links}
\author{John Archibald Wheeler}
\date{}

\begin{document}
	\maketitle
	
	\begin{center}
		\small Reproduced from \emph{Proc. 3rd Int. Symp. Foundations of Quantum Mechanics}, Tokyo, 1989, pp. 354--368.\\
		\small Copyright of this paper is held by the author.
	\end{center}
	
	\begin{abstract}
		This report reviews what quantum physics and information theory have to tell us
		about the age-old question, How come existence? No escape is evident from four
		conclusions: (1) The world cannot be a giant machine, ruled by any preestablished
		continuum physical law. (2) There is no such thing at the microscopic level as
		space or time or spacetime continuum. (3) The familiar probability function or
		functional, and wave equation or functional wave equation, of standard quantum
		theory provide mere continuum idealizations and by reason of this circumstance
		conceal the information-theoretic source from which they derive. (4) No element in
		the description of physics shows itself as closer to primordial than the elementary
		quantum phenomenon, that is, the elementary device-intermediated act of posing a
		yes-no physical question and eliciting an answer or, in brief, the elementary act of
		observer-participancy. Otherwise stated, every physical quantity, every it, derives
		its ultimate significance from bits, binary yes-or-no indications, a conclusion which
		we epitomize in the phrase, it from bit.
	\end{abstract}
	
	\section{Quantum Physics Requires a New View of Reality}
	Revolution in outlook though Kepler, Newton, and Einstein brought us [1--4], and
	still more startling the story of life [5--7] that evolution forced upon an unwilling
	world, the ultimate shock to preconceived ideas lies ahead, be it a decade hence,
	a century or a millennium. The overarching principle of 20th-century physics, the
	quantum [8]---and the principle of complementarity [9] that is central idea of the
	quantum---leaves us no escape, Niels Bohr tells us, [10] from ``a radical revision
	of our attitude as regards physical reality'' and a ``fundamental modification of
	all ideas regarding the absolute character of physical phenomena.'' Transcending
	Einstein's summons [11] of 1908, ``This quantum business is so incredibly important
	and difficult that everyone should busy himself with it,'' Bohr's modest words direct
	us to the supreme goal: Deduce the quantum from an understanding of existence.
	
	How make headway toward a goal so great against difficulties so large? The
	search for understanding presents to us three questions, four no's and five clues:
	
	Three questions,
	\begin{itemize}[leftmargin=2em]
		\item How come existence?
		\item How come the quantum?
		\item How come ``one world'' out of many observer-participants?
	\end{itemize}
	
	Four no's,
	\begin{itemize}[leftmargin=2em]
		\item No tower of turtles
		\item No laws
		\item No continuum
		\item No space, no time.
	\end{itemize}
	
	Five clues,
	\begin{itemize}[leftmargin=2em]
		\item The boundary of a boundary is zero.
		\item No question? No answer!
		\item The super-Copernican principle.
		\item ``Consciousness''
		\item More is different.
	\end{itemize}
	
	\section{``It from Bit'' as Guide in Search for Link Connecting Physics, Quantum and Information}
	In default of a tentative idea or working hypothesis, these questions, no's and clues
	--- yet to be discussed --- do not move us ahead. Nor will any abundance of clues
	assist a detective who is unwilling to theorize how the crime was committed! A
	wrong theory? The policy of the engine inventor, John Kris, reassures us, ``Start
	her up and see why she don't go!'' In this spirit [12--47] I, like other searchers [48--51]
	attempt formulation after formulation of the central issues, and here present a wider
	overview, taking for working hypothesis the most effective one that has survived this
	winnowing: It from bit. Otherwise put, every it---every particle, every field of
	force, even the spacetime continuum itself---derives its function, its meaning, its
	very existence entirely---even if in some contexts indirectly---from the apparatus-
	elicited answers to yes or no questions, binary choices [52], bits.
	
	It from bit symbolizes the idea that every item of the physical world has at
	bottom---at a very deep bottom, in most instances---an immaterial source and
	explanation; that what we call reality arises in the last analysis from the posing of
	yes-no questions and the registering of equipment-evoked responses; in short, that
	all things physical are information-theoretic in origin and this is a participatory
	universe.
	
	Three examples may illustrate the theme of it from bit. First, the photon. With
	polarizer over the distant source and analyzer of polarization over the photodetector
	under watch, we ask the yes or no question, ``Did the counter register a click during
	the specified second?'' If yes, we often say, ``A photon did it.'' We know perfectly
	well that the photon existed neither before the emission nor after the detection.
	However, we also have to recognize that any talk of the photon ``existing'' during
	the intermediate period is only a blown-up version of the raw fact, a count.
	
	The yes or no that is recorded constitutes an unsplittable bit of information. A
	photon, Wootters and Zurek demonstrate [53, 54], cannot be cloned.
	
	As second example of it from bit, we recall the Aharonov-Bohm scheme [55] to
	measure a magnetic flux. Electron counters stationed off to the right of a doubly-
	slit screen give yes-or-no indications of the arrival of an electron from the source
	located off to the left of the screen, both before the flux is turned on and afterward.
	That flux of magnetic lines of force finds itself embraced between---but untouched
	by---the two electron beams that fan out from the two slits. The beams interfere.
	The shift in interference fringes between field off and field on reveals the magnitude
	of the flux,
	
	\begin{align*}
		(\text{phase change around perimeter of the included area}) &= 2\pi \times (\text{shift of interference pattern, measured in number of fringes}) \\
		&= (\text{electron charge}) \times (\text{magnetic flux embraced}) / \hbar c
		\tag{19.1}
	\end{align*}
	
	Here $\hbar = 1.0546 \times 10^{-27} \text{g cm}^2/\text{s}$ is the quantum in conventional units, or in geometric units [4, 16]---where both time and mass are measured in the units of length---$\hbar = \hbar c = 2.612 \times 10^{-66} \text{cm}^2 =$ the square of the Planck length, $1.616 \times 10^{-33} \text{cm} =$ what we hereafter term the Planck area.
	
	Not only in electrodynamics but also in geometrodynamics and in every other
	gauge-field theory, as Anandan, Aharonov and others point out [56, 57] the difference around a circuit in the phase of an appropriately chosen quantum-mechanical
	probability amplitude provides a measure of the field. Here again the concept of it
	from bit applies [38]. Field strength or spacetime curvature reveals itself through
	shift of interference fringes, fringes that stand for nothing but a statistical pattern
	of yes-or-no registrations.
	
	When a magnetometer reads that it which we call a magnetic field, no reference
	at all to a bit seems to show itself. Therefore we look closer. The idea behind the
	operation of the instrument is simple. A wire of length $l$ carries a current $i$ through
	
	\begin{figure}[htbp]
		\centering
		\fbox{\parbox{0.78\linewidth}{\centering Figure 19.1 from the source PDF.}}
		\caption{Symbolic representation of the ``telephone number'' of the particular one of the
			$2^N$ conceivable, but by now indistinguishable, configurations out of which this particular
			black hole, of Bekenstein number $N$ and horizon area $4N\hbar \ln 2$, was put together. Symbol,
			also, in a broader sense, of the theme that every physical entity, every it, derives from bits.
			Reproduced from JGST, p.220.}
	\end{figure}
	
	a magnetic field $B$ that runs perpendicular to it. In consequence the piece of copper
	receives in the time $t$ a transfer of momentum $p$ in a direction $z$ perpendicular to
	the directions of the wire and of the field,
	
	\begin{align*}
		p &= B l i t \\
		&= (\text{flux per unit } z) \times (\text{charge, } e, \text{ of the elementary carrier of current}) \\
		&\quad \times (\text{number, } N, \text{ of carriers that pass in the time } t)
		\tag{19.2}
	\end{align*}
	
	This impulse is the source of the force that displaces the indicator needle of the
	magnetometer and gives us an instrument reading. We deal with bits wholesale
	rather than bits retail when we run the fiducial current through the magnetometer
	coil, but the definition of field founds itself no less decisively on bits.
	
	As third and final example of it from bit we recall the wonderful quantum
	finding of Bekenstein [58--60]---totally unexpected denouement of earlier classical
	work of Penrose [61] Christodoulou [62] and Ruffini [63]---refined by Hawking [64,
	65] that the surface area of the horizon of a black hole, rotating or not, measures
	the entropy of the black hole. Thus this surface area, partitioned in imagination
	(Fig. 19.1) into domains each of size $4\hbar \ln 2$, that is, $2.77...$ times the Planck area,
	yields the Bekenstein number, $N$; and the Bekenstein number, so Thorne and Zurek
	explain [66] tells us the number of binary digits, the number of bits, that would be
	required to specify in all detail the configuration of the constituents out of which
	the black hole was put together. Entropy is a measure of lost information. To no
	community of newborn outside observers can the black hole be made to reveal out
	of which particular one of $2^N$ configurations it was put together. Its size, an it, is
	fixed by the number, $N$, of bits of information hidden within it.
	
	The quantum, $\hbar$, in whatever correct physics formula it appears, thus serves as
	lamp. It lets us see horizon area as information lost, understand wave number of
	light as photon momentum and think of field flux as bit-registered fringe shift.
	Giving us its as bits, the quantum presents us with physics as information.
	
	How come a value for the quantum so small as $\hbar = 2.612 \times 10^{-66} \text{cm}^2$? As well
	as ask why the speed of light is so great as $c = 3 \times 10^{10} \text{cm/s}$! No such constant
	as the speed of light ever makes an appearance in a truly fundamental account
	of special relativity or Einstein geometrodynamics, and for a simple reason: Time
	and space are both tools to measure interval. We only then properly conceive
	them when we measure them in the same units [4, 16]. The numerical value of
	the ratio between the second and the centimeter totally lacks teaching power. It
	is an historical accident. Its occurrence in equations obscured for decades one of
	Nature's great simplicities. Likewise with $\hbar$; every equation that contains an $\hbar$
	floats a banner, ``It from bit''. The formula displays a piece of physics that we
	have learned to translate into information-theoretic terms. Tomorrow we will have
	learned to understand and express all of physics in the language of information. At
	that point we will revalue $\hbar = 2.612 \times 10^{-66} \text{cm}^2$---as we downgrade $c = 3 \times 10^{10} \text{cm/s}$ today---from constant of Nature to artifact of history, and from foundation
	of truth to enemy of understanding.
	
	\section{Four No's}
	To the question, ``How come the quantum?'' we thus answer, ``Because what we
	call existence is an information-theoretic entity.'' But how come existence? Its
	as bits, yes; and physics as information, yes; but whose information? How does
	the vision of one world arise out of the information-gathering activities of many
	observer-participants? In the consideration of these issues we adopt for guidelines
	four no's.
	
	First no: ``No tower of turtles,'' advised William James. Existence is not a globe
	supported by an elephant, supported by a turtle, supported by yet another turtle,
	and so on. In other words, no infinite regress. No structure, no plan of organization,
	no framework of ideas underlaid by another structure or level of ideas, underlaid
	by yet another level, by yet another, ad infinitum, down to a bottomless night. To
	endlessness no alternative is evident but loop [47, 67], such a loop as this: Physics
	gives rise to observer-participancy; observer-participancy gives rise to information;
	and information gives rise to physics.
	
	Existence thus built [68] on ``insubstantial nothingness''? Rutherford and Bohr
	made a table no less solid when they told us it was $99.9...$ percent emptiness.
	Thomas Mann may exaggerate when he suggests [69] that ``\ldots we are actually
	bringing about what seems to be happening to us,'' but Leibniz [70] reassures us,
	``Although the whole of this life were said to be nothing but a dream and the physical
	world nothing but a phantasm, I should call this dream or phantasm real enough
	if, using reason well, we were never deceived by it.''
	
	Second no: No laws. ``So far as we can see today, the laws of physics cannot
	have existed from everlasting to everlasting. They must have come into being at
	the big bang. There were no gears and pinions, no Swiss watch-makers to put
	things together, not even a pre-existing plan\ldots Only a principle of organization
	which is no organization at all would seem to offer itself. In all of mathematics,
	nothing of this kind more obviously offers itself than the principle that 'the boundary
	of boundary is zero.' Moreover, all three great field theories of physics use this
	principle twice over\ldots This circumstance would seem to give us some reassurance
	that we are talking sense when we think of\ldots physics being'' [32] as foundation-free
	as a logic loop, the closed circuit of ideas in a self-referential deductive axiomatic
	system [71--74].
	
	Universe as machine? This universe one among a great ensemble of machine
	universes, each differing from the others in the values of the dimensionless constants
	of physics? Our own selected from this ensemble by an anthropic principle of one
	or another form [75]? We reject here the concept of universe as machine not least
	because it ``has to postulate explicitly or implicitly, a supermachine, a scheme,
	a device, a miracle, which will turn out universes in infinite variety and infinite
	number'' [47].
	
	Directly opposite to the concept of universe as machine built on law is the vision
	of a world self-synthesized. On this view, the notes struck out on a piano by the
	observer-participants of all places and all times, bits though they are, in and by
	themselves constitute the great wide world of space and time and things.
	
	Third no: No continuum. No continuum in mathematics and therefore no continuum in physics. A half-century of development in the sphere of mathematical
	logic [76] has made it clear that there is no evidence supporting the belief in the existential character of the number continuum. ``Belief in this transcendental world,''
	Hermann Weyl tells us, ``taxes the strength of our faith hardly less than the doctrines of the early Fathers of the Church or of the scholastic philosophers of
	the Middle Ages'' [77]. This lesson out of mathematics applies with equal strength
	to physics. ``Just as the introduction of the irrational numbers\ldots is a convenient
	myth [which] simplifies the laws of arithmetic\ldots so physical objects,'' Willard Van
	Orman Quine tells us [78] ``are postulated entities which round out and simplify
	our account of the flux of existence\ldots The conceptual scheme of physical objects
	is a convenient myth, simpler than the literal truth and yet containing that literal
	truth as a scattered part.''
	
	Nothing so much distinguishes physics as conceived today from mathematics as
	the difference between the continuum character of the one and the discrete character of the other. Nothing does so much to extinguish this gap as the elementary
	quantum phenomenon ``brought to a close,'' as Bohr puts it [10] by ``an irreversible
	act of amplification,'' such as the click of a photodetector or the blackening of
	a grain of photographic emulsion. Irreversible? More than one idealized experiment [38] illustrates how hard it is, even today, to give an all-inclusive definition of
	the term irreversible. Those difficulties supply pressure, however, not to retreat to
	old ground, but to advance to new insight. In brief, continuum-based physics, no;
	information-based physics, yes.
	
	Fourth and last no: No space, no time. Heaven did not hand down the word
	``time''. Man invented it, perhaps positing hopefully as he did that ``Time is Nature's
	way to keep everything from happening all at once'' [79]. If there are problems with
	the concept of time, they are of our own creation! As Leibniz tells us, [80] ``\ldots time
	and space are not things, but orders of things\ldots ;'' or as Einstein put it, [81] ``Time
	and space are modes by which we think, and not conditions in which we live.''
	
	What are we to say about that weld of space and time into spacetime which
	Einstein gave us in his 1915 and still standard classical geometrodynamics? On this
	geometry quantum theory, we know, imposes fluctuations [13, 14, 82]. Moreover, the
	predicted fluctuations grow so great at distances of the order of the Planck length
	that in that domain they put into question the connectivity of space and deprive
	the very concepts of ``before'' and ``after'' of all meaning [83]. This circumstance
	reminds us anew that no account of existence can ever hope to rate as fundamental
	which does not translate all of continuum physics into the language of bits.
	
	We will not feed time into any deep-reaching account of existence. We must
	derive time---and time only in the continuum idealization---out of it. Likewise
	with space.
	
	\section{Five Clues}
	First clue: The boundary of a boundary is zero. This central principle of algebraic
	topology [84], identity, triviality, tautology, though it is, is also the unifying theme
	of Maxwell electrodynamics, Einstein geometrodynamics and almost every version
	of modern field theory [42], [85--88]. That one can get so much from so little, almost
	everything from almost nothing, inspires hope that we will someday complete the
	mathematization of physics and derive everything from nothing, all law from no
	law.
	
	Second clue: No question, no answer. Better put, no bit-level question, no bit-level answer. So it is in the game of twenty questions in its surprise version [89].
	And so it is for the electron circulating within the atom or a field within a space.
	To neither field nor particle can we attribute a coordinate or momentum until a
	device operates to measure the one or the other. Moreover any apparatus that
	accurately [90] measures the one quantity inescapably rules out then and there the
	operation of equipment to measure the other [9, 91, 92]. In brief, the choice of
	question asked, and choice of when it's asked, play a part---not the whole part,
	but a part---in deciding what we have the right to say [38, 43].
	
	Bit-registration of a chosen property of the electron, a bit-registration of the
	arrival of a photon, Aharonov-Bohm bit-based determination of the magnitude of a
	field flux, bulk-based count of bits bound in a black hole: All are examples of physics
	expressed in the language of information. However, into a bit count that one might
	have thought to be a private matter, the rest of the nearby world irresistibly thrusts
	itself. Thus the atom-to-atom distance in a ruler---basis for a bit count of distance---
	evidently has no invariant status, depending as it does on the temperature and
	pressure of the environment. Likewise the shift of fringes in the Aharonov-Bohm
	experiment depends not only upon the magnetic flux itself, but also on the charge
	of the electron. But this electron charge---when we take the quantum itself to
	be Nature's fundamental measuring unit---is governed by the square root of the
	quantity $e^2 / \hbar c \approx 1/137.036...$, a ``constant'' which---for extreme conditions---
	is as dependent on the local environment [93] as is a dielectric ``constant'' or the
	atom-to-atom spacing in the ruler.
	
	The contribution of the environment becomes overwhelmingly evident when we
	turn from length of bar or flux of field to the motion of alpha particle through cloud
	chamber, dust particle through $3$ K background radiation or Moon through space.
	This we know from the analyses of Bohr and Mott [94], Zeh [95, 96], Joos and
	Zeh [97], Zurek [98--100] and Unruh and Zurek [101]. It from bit, yes; but the rest
	of the world also makes a contribution, a contribution that suitable experimental
	design can minimize but not eliminate. Unimportant nuisance? No. Evidence the
	whole show is wired up together? Yes. Objection to the concept of every it from
	bits? No.
	
	Build physics, with its false face of continuity, on bits of information! What
	this enterprise is we perhaps see more clearly when we examine for a moment a
	thoughtful, careful, wide-reaching exposition [102] of the directly opposite thesis,
	that physics at bottom is continuous; that the bit of information is not the basic entity. Rate as false the claim that the bit of information is the basic entity. Instead,
	attempt to build everything on the foundation of some ``grand unified field theory''
	such as string theory [103, 104]---or in default of that, on Einstein's 1915 and still
	standard geometrodynamics. Hope to derive that theory by way of one or another
	plausible line of reasoning. But don't try to derive quantum theory. Treat it as
	supplied free of charge from on high. Treat quantum theory as a magic sausage
	grinder which takes in as raw meat this theory, that theory or the other theory
	and turns out a ``wave equation,'' one solution of which is ``the'' wave function for
	the universe [14, 102, 105--107]. From start to finish accept continuity as right and
	natural: Continuity in the manifold, continuity in the wave equation, continuity in
	its solution, continuity in the features that it predicts. Among conceivable solutions
	of this wave equation select as reasonable one which ``maximally decoheres,'' one
	which exhibits ``maximal classicity''---maximal classicity by reason, not of ``something external to the framework of wave function and Schrodinger equation,'' but
	something in ``the initial conditions of the universe specified within quantum theory
	itself.''
	
	How compare the opposite outlooks of decoherence and it-from-bit? Remove
	the casing that surrounds the workings of a giant computer. Examine the bundles
	of wires that run here and there. What is the status of an individual wire? Mathematical limit of bundle? Or building block of bundle? The one outlook regards
	the wave equation and wave function to be primordial and precise and built on
	continuity, and the bit to be idealization. The other outlook regards the bit to be
	the primordial entity, and wave equation and wave function to be secondary and
	approximate---and derived from bits via information theory.
	
	Derived, yes; but how? No one has done more than William Wootters towards
	opening up a pathway [108, 109] from information to quantum theory. He puts
	into connection two findings, long known, but little known. Already before the
	advent of wave mechanics, he notes, the analyst of population statistics R. A. Fisher
	proved [110, 111] that the proper tool to distinguish one population from another is
	not the probability of this gene, that gene and the third gene (for example), but the
	square roots of these probabilities; that is to say, the two probability amplitudes,
	each probability amplitude being a vector with three components. More precisely,
	Wootters proves, the distinguishability between the two populations is measured by
	the angle in Hilbert space between the two state vectors, both real. Fisher, however,
	was dealing with information that sits ``out there''. In microphysics, however, the
	information does not sit out there. Instead, Nature in the small confronts us with
	a revolutionary pistol, ``No question, no answer.'' Complementarity rules. And
	complementarity as E.C.G. Stueckelberg proved [112, 113] as long ago as 1952,
	and as Saxon made more readily understandable [114] in 1964, demands that the
	probability amplitudes of quantum physics must be complex. Thus Wootters derives
	familiar Hilbert space with its familiar complex probability amplitudes from the twin
	demands of complementarity and measure of distinguishability.
	
	Try to go on from Wootters's finding to deduce the full blown machinery of
	quantum field theory? Exactly not to try to do so---except as idealization---is
	the demand laid on us by the concept of it from bit. How come?
	
	Probabilities exist ``out there'' no more than do space or time or the position of
	the atomic electron. Probability, like time, is a concept invented by humans, and
	humans have to bear the responsibility for the obscurities that attend it. Obscurities
	there are whether we consider probability defined as frequency [115] or defined a la
	Bayes [116--119]. Probability in the sense of frequency has no meaning as applied
	to the spontaneous fission of the particular plutonium nucleus that triggered the
	November 1, 1952 H-bomb blast.
	
	What about probabilities of a Bayesian cast, probabilities ``interpreted not as frequencies observable through experiments, but as degrees of plausibility one assigns
	to each hypothesis based on the data and on one's assessment of the plausibility
	of the hypotheses prior to seeing the data'' [120]. Belief-dependent probabilities,
	different probabilities assigned to the same proposition by different people [121]?
	Probabilities associated [122] with the view that ``objective reality is simply an
	interpretation of data agreed to by large numbers of people?''
	
	Heisenberg directs us to the experiences [123] of the early nuclear-reaction-rate
	theorist Fritz Houtermans, imprisoned in Kharkov during the time of the Stalin
	terror, ``\ldots the whole cell would get together to produce an adequate confession
	\ldots [and] helped them [the prisoners] to compose their 'legends' and phrase them
	properly, implicating as few others as possible.''
	
	Existence as confession? Myopic but in some ways illuminating formulation of
	the demand for intercommunication implicit in the theme of it from bit!
	
	So much for ``No question, no answer.''
	
	Third clue: The super-Copernican principle [47]. This principle rejects now-
	centeredness in any account of existence as firmly as Copernicus repudiated here-
	centeredness. It repudiates most of all any tacit adoption of here-centeredness in
	assessing observer-participants and their number.
	
	What is an observer-participant? One who operates an observing device and
	participates in the making of meaning, meaning in the sense of Føllesdal [124],
	``Meaning is the joint product of all the evidence that is available to those who
	communicate.'' Evidence that is available? The investigator slices a rock and photographs the evidence for the heavy nucleus that arrived in the cosmic radiation of a
	billion years ago [38]. Before he can communicate his findings, however, an asteroid
	atomizes his laboratory, his records, his rocks and him. No contribution to meaning!
	Or at least no contribution then. A forensic investigation of sufficient detail and
	wit to reconstruct the evidence of the arrival of that nucleus is difficult to imagine.
	
	What about the famous tree that fell in the forest with no one around [125]? It
	leaves a fallout of physical evidence so near at hand and so rich that a team of up-
	to-date investigators can establish what happened beyond all doubt. Their findings
	contribute to the establishment of meaning.
	
	``Measurements and observations,'' it has been said, [102] ``cannot be fundamental notions in a theory which seeks to discuss the early universe when neither
	existed.'' On this view the past has a status beyond all questions of observer-
	participancy. It from bit offers us a different vision: ``Reality is theory'' [126]; ``the
	past has no evidence except as it is recorded in the present'' [127]. The photon that
	we are going to register tonight from that four billion-year old quasar cannot be said
	to have had an existence ``out there'' three billion years ago, or two (when it passed
	an intervening gravitational lens) or one, or even a day ago. Not until we have fixed
	arrangements at our telescope do we register tonight's quantum as having passed
	to the left (or right) of the lens or by both routes (as in a double slit experiment).
	This registration like every delayed-choice experiment [21, 40], reminds us that no
	elementary quantum phenomenon is a phenomenon until, in Bohr's words [10], ``It
	has been brought to a close'' by ``an irreversible act of amplification.'' What we call
	the past is built on bits.
	
	Enough bits to structure a universe so rich in features as we know this world to
	be. Preposterous! Mice and men and all on Earth who may ever come to rank as
	intercommunicating meaning-establishing observer-participants will never mount a
	bit count sufficient to bear so great a burden.
	
	The count of bits needed, huge though it may be, nevertheless, so far as we
	can judge, does not reach infinity. In default of a better estimate, we follow familiar reasoning [128] and translate into the language of the bits the entropy of
	the primordial cosmic fireball as deduced from the entropy of the present $2.735$ K
	(uncertainty $<0.05$ K) microwave relict radiation [129] totaled over a 3-sphere of
	radius $13.2 \times 10^9$ light years (uncertainty $<35\%$) [130] or $1.25 \times 10^{28}$ cm and of
	volume $2\pi^2 \text{radius}^3$,
	
	\begin{align*}
		(\text{number of bits}) &= (\log_2 e) \times (\text{number of nats}) \\
		&= (\log_2 e) \times (\text{entropy / Boltzmann's constant, } k) \\
		&= 1.44... \times [(8\pi^4/45)(\text{radius} \cdot kT/\hbar c)^3] \\
		&= 8 \times 10^{88}
		\tag{19.3}
	\end{align*}
	
	It would be totally out of place to compare this overpowering number with the number of bits of information elicited up to date by observer-participancy. So warns
	the super-Copernican principle. We today, to be sure, through our registering devices, give a tangible meaning to the history of the photon that started on its way
	from a distant quasar long before there was any observer-participancy anywhere.
	However, the far more numerous establishers of meaning of time to come have a
	like inescapable part---by device-elicited questions and registration of answer---
	in generating the ``reality'' of today. For this purpose, moreover, there are billions of years yet to come, billions on billions of sites of observer-participancy yet
	to be occupied. How far foot and ferry have carried meaning-making communication in fifty thousand years gives faint feel for how far interstellar propagation is
	destined [131, 132] to carry it in fifty billion years.
	
	Do bits needed balance bits achievable? They must, declares the concept of
	``world as system self-synthesized by quantum networking'' [47]. By no prediction
	does this concept more clearly expose itself to destruction, in the sense of Popper [133].
	
	Fourth clue: ``Consciousness''. We have traveled what may seem a dizzying path.
	First, elementary quantum phenomenon brought to a close by an irreversible act
	of amplification. Second, the resulting information expressed in the form of bits.
	Third, this information used by observer-participants---via communication---to
	establish meaning. Fourth, from the past through the billenniums to come, so many
	observer-participants, so many bits, so much exchange of information, as to build
	what we call existence.
	
	Doesn't this it-from-bit view of existence seek to elucidate the physical world,
	about which we know something, in terms of an entity about which we know almost
	nothing, consciousness [134--137]? And doesn't Marie Sklodowska Curie tell us,
	``Physics deals with things, not people?'' Using such and such equipment, making
	such and such a measurement, I get such and such a number. Who I am has nothing
	to do with this finding. Or does it? Am I sleepwalking [138, 139]? Or am I one
	of those poor souls without the critical power to save himself from pathological
	science [140--142]?
	
	Under such circumstances any claim to have ``measured'' something falls flat
	until it can be checked out with one's fellows. Checked how? Morton White reminds us [143] how the community applies its tests of credibility, and in this connection quotes analyses by Chauncey Wright, Josiah Royce and Charles Saunders
	Peirce [144]. Parmenides of Elea [145] (c. 515 B.C.--450+ B.C.) may tell us that
	``What is\ldots is identical with the thought that recognizes it.'' We, however, steer
	clear of the issues connected with ``consciousness.'' The line between the unconscious and the conscious begins to fade [146] in our day as computers evolve and
	develop---as mathematics has---level upon level upon level of logical structure.
	We may someday have to enlarge the scope of what we mean by a ``who''. This
	granted, we continue to accept---as essential part of the concept of it from bit---
	Føllesdal's guideline [124], ``Meaning is the joint product of all the evidence that is
	available to those who communicate.'' What shall we say of a view of existence [147]
	that appears, if not anthropomorphic in its use of the word ``who,'' still overly centered on life and consciousness? It would seem more reasonable to dismiss for
	the present the semantic overtones of ``who'' and explore and exploit the insights to
	be won from the phrases, ``communication'' and ``communication employed to establish meaning.''
	
	Føllesdal's statement supplies, not an answer, but the doorway to new questions.
	For example, man has not yet learned how to communicate with ant. When he
	does, will the questions put to the world around by the ant and the answers that he
	elicits contribute their share, too, to the establishment of meaning? As another issue
	associated with communication, we have yet to learn how to draw the line between
	a communication network that is closed, or parochial, and one that is open. And
	how to use that difference to distinguish between reality and poker---or another
	game [148, 149]---so intense as to appear more real than reality. No term in
	Føllesdal's statement poses greater challenge to reflection than ``communication,''
	descriptor of a domain of investigation [150--152] that enlarges in sophistication with
	each passing year.
	
	Fifth and final clue: More is different [153]. Not by plan but by inner necessity
	a sufficiently large number of H$_2$O molecules collected in a box will manifest solid,
	liquid and gas phases. Phase changes, superfluidity and superconductivity all bear
	witness to Anderson's pithy point, more is different.
	
	We do not have to turn to objects so material as electrons, atoms and molecules
	to see big numbers generating new features. The evolution from small to large has
	already in a few decades forced on the computer a structure [154, 155] reminiscent
	of biology by reason of its segregation of different activities into distinct organs.
	Distinct organs, too, the giant telecommunications system of today finds itself inescapably evolving [151, 152]. Will we someday understand time and space and all
	the other features that distinguish physics---and existence itself---as the similarly
	self-generated organs of a self-synthesized information system [156--158]?
	
	\section{Conclusion}
	The spacetime continuum? Even continuum existence itself? Except as idealization neither the one entity nor the other can make any claim to be a primordial
	category in the description of Nature. It is wrong, moreover, to regard this or that
	physical quantity as sitting ``out there'' with this or that numerical value in default
	of question asked and answer obtained by way of appropriate observing device.
	The information thus solicited makes physics and comes in bits. The count of bits
	drowned in the dark night of a black hole displays itself as horizon area, expressed
	in the language of Bekenstein number. The bit count of the cosmos, however it is
	figured, is ten raised to a very large power. So also is the number of elementary
	acts of observer-participancy over any time of the order of fifty billion years. And,
	except via those time-leaping quantum phenomena that we rate as elementary acts
	of observer-participancy, no way has ever offered itself to construct what we call
	``reality.'' That's why we take seriously the theme of it from bit.
	
	\section{Agenda}
	Intimidating though the problem of existence continues to be, the theme of it from
	bit breaks it down into six issues that invite exploration:
	
	One: Go beyond Wootters and determine what, if anything, has to be added to
	distinguishability and complementarity to obtain all of standard quantum theory.
	
	Two: Translate the quantum versions of string theory and of Einstein's geometrodynamics from the language of continuum to the language of bits.
	
	Three: Sharpen the concept of bit. Determine whether ``an elementary quantum
	phenomenon brought to a close by an irreversible act of amplification'' has at bottom
	(1) the 0-or-1 sharpness of definition of bit number nineteen in a string of binary
	digits, or (2) the accordion property of a mathematical theorem, the length of which,
	that is, the number of supplementary lemmas contained in which, the analyst can
	stretch or shrink according to his convenience.
	
	Four: Survey one by one with an imaginative eye the powerful tools that mathematics---including mathematical logic---has won and now offers to deal with
	theorems on a wholesale rather than a retail level, and for each such technique
	work out the transcription into the world of bits. Give special attention to one and
	another deductive axiomatic system which is able to refer to itself [159], one and
	another self-referential deductive system.
	
	Five: From the wheels-upon-wheels-upon-wheels evolution of computer programming dig out, systematize and display every feature that illuminates the level-
	upon-level-upon-level structure of physics.
	
	Six: Capitalize on the findings and outlooks of information theory [160--163],
	algorithmic entropy [164], evolution of organisms [165--167] and pattern recognition [168--175]. Search out every link each has with physics at the quantum level.
	Consider, for instance, the string of bits 1111111\ldots and its representation as the
	sum of the two strings 1001110\ldots and 0110001\ldots Explore and exploit the connection between this information-theoretic statement and the findings of theory and
	experiment on the correlation between the polarizations of the two photons emitted
	in the annihilation of singlet positronium [176] and in like Einstein-Podolsky-Rosen
	experiments [177]. Seek out, moreover, every realization in the realm of physics of
	the information-theoretic triangle inequality recently discovered by Zurek [178].
	
	Finally: Deplore? No, celebrate the absence of a clean clear definition of the
	term ``bit'' as elementary unit in the establishment of meaning. We reject ``that
	view of science which used to say, 'Define your terms before you proceed.' The
	truly creative nature of any forward step in human knowledge,'' we know, ``is such
	that theory, concept, law and method of measurement---forever inseparable---are
	born into the world in union [179].'' If and when we learn how to combine bits in
	fantastically large numbers to obtain what we call existence, we will know better
	what we mean both by bit and by existence.
	
	A single question animates this report: Can we ever expect to understand existence? Clues we have, and work to do, to make headway on that issue. Surely
	someday, we can believe, we will grasp the central idea of it all as so simple, so
	beautiful, so compelling that we will all say to each other, ``Oh, how could it have
	been otherwise! How could we all have been so blind so long!'
	
	\section*{Acknowledgements}
	For discussion, advice or judgment on one or another issue taken up in this review,
	I am indebted to Nandar Balazs, John D. Barrow, Charles H. Bennett, David
	Deutsch, Robert H. Dicke, Freeman Dyson and the late Richard P. Feynman as
	well as David Gross, James B. Hartle, John J. Hopfield, Paul C. Jeffries, Bernulf
	Kanitscheider, Arkady Kheyfets and Rolf W. Landauer; and to Warner A. Miller,
	John R. Pierce, Willard Van Orman Quine, Benjamin Schumacher and Frank J.
	Tipler as well as William G. Unruh, Morton White, Eugene P. Wigner, William K.
	Wootters, Hans Dieter Zeh and Wojciech H. Zurek. For assistance in preparation
	of this report I thank E. L. Bennett and NSF grant PHY245--6243 to Princeton
	University. I give special thanks to the sponsors of the 28--31 August 1989 conference
	ISQM Tokyo '89 at which the then current version of the present analysis was
	reported.
	
	This report evolved from presentations at Santa Fe Institute Conferences, 29
	May--2 June and 4--8 June 1989 and at the 3rd International Symposium on Foundations of Quantum Mechanics in the Light of New Technology, Tokyo, 28--31 August
	1989, under the title ``Information, Physics, Quantum: The Search for Links;'' and
	headed ``Can We Ever Expect to Understand Existence?'', as the Penrose Lecture
	at the 20--22 April 1989 annual meeting of Benjamin Franklin's ``American Philosophical Society, held at Philadelphia for Promoting Useful Knowledge,'' and the
	Accademia Nazionale dei Lincei Conference on La Verita nella Scienza, Rome, 13
	October 1989; submitted to proceedings of all four in fulfillment of obligation and
	in deep appreciation for hospitality. Preparation for publication assisted in part by
	NSF Grant PHY 245--6243 to Princeton University.
	
	\section*{Discussion}
	A discussion followed:
	
	\noindent\textbf{N. G. van Kampen:} Did you mean to say that the observer influences the observed object?
	
	\noindent\textbf{J. A. Wheeler:} The observer does not influence the past. Instead, by his choice of question, he decides about what feature of the object he shall have the right to make a clear statement.
	
	\noindent\textbf{J.P. Vigier:} Two problems.
	1. The first is that the QSO raise lots of unsolved problems, i.e.---strange quantized $N/\log(1+z)$ relation---correlation with galaxies (Arp)---angular correlation with brightest nearby galaxies (Burbidge et al.)
	2. The second is that the idea (Einstein et al.) of the reality of fields has led (assuming that ``particles'' are field singularities) to the only known justification of the geodesic law. To contest it is to make the meaning of dynamical behaviour purely observer-dependent, i.e., to kill the reality of the physical world.
	
	\noindent\textbf{J.A. Wheeler:}
	1. The book by Thorne and colleagues, ``Black Holes: The Membrane Paradigm,'' describes how a supermassive black hole, endowed via accretion with great angular momentum inside and an accretion disk outside, produces counter-directed jets and radiation of great power. I know no other mechanism able to produce quasars.
	2. No one has discovered a way to get a particle of wave length $\lambda$ from point A through empty flat space to a point B at a great distance $L$ without its undergoing on the way a transverse spread of the order $\sqrt{L\lambda}$. This spread imposes an inescapable limitation on the classical concept of ``worldline.''
	
	\section*{References}
	
	\begin{thebibliography}{180}
		
		\bibitem{1}
		J. Kepler (1571--1630): \emph{Harmonices Mundi}, 5 books (1619). The appendix to Kepler's Book 5 contains one side, the publications of the English physician and thinker Robert Fludd (1574--1637) the other side, of a great debate, analysed by Wolfgang Pauli [W. Pauli: ``Der Einfluss archetypischer Vorstellungen auf die Bildung naturwissenschaftlicher Theorien bei Kepler'' in \emph{Naturerkl\"arung und Psyche} (Rascher, Z\"urich, 1952) p.109--194; reprinted in \emph{Wolfgang Pauli: Collected Scientific Papers}, eds. R. Kronig and V. F. Weisskopf (Interscience-Wiley, New York, 1964) Vol.1, p.1023]. Totally in contrast to Fludd's concept of intervention from on high [\emph{Utriusque Cosmo Maioris scilicet et Minoris Metaphysica, Physica atque technica Historia}, 1st ed. (Oppenheim, 1621)] was Kepler's guiding principle, \emph{Ubi materia, ibi geometria}---where there is matter, there is geometry. It was not directly from Kepler's writings, however, that Newton learned of Kepler's three great geometry-driven findings about the motions of the planets in space and in time, but from the distillation of Kepler offered by Thomas Streete (1622--1689), \emph{Astronomia Carolina: A New Theorie of the Celestial Motions} (London, 1661).
		
		\bibitem{2}
		I. Newton: \emph{Philosophiae naturalis principia mathematica}, 1st ed. (London, 1687).
		
		\bibitem{3}
		A. Einstein: ``Zur allgemeinen Relativit\"atstheorie'' \emph{Preuss. Akad. Wiss. Berlin, Sitzber} (1915) p.799--801; also (1915) p. 832--839, 844--847; (1916) p.688--696 and (1917) p.142--152.
		
		\bibitem{4}
		J. A. Wheeler: \emph{Journey into Gravity and Spacetime} (Scientific American Library, Freeman, New York, 1990), cited hereafter as JGST.
		
		\bibitem{5}
		J. G. Mendel: ``Versuche \"uber Pflanzenhybriden'' \emph{Verhandlungen des Naturforschenden Vereins in Br\"unn} 4 (1866).
		
		\bibitem{6}
		C. R. Darwin (1809--1882): \emph{On the Origin of Species by Means of Natural Selection, or the Preservation of Favoured Races in the Struggle for Life} (London, 1859).
		
		\bibitem{7}
		J. D. Watson and F. H. C. Crick: ``Molecular structure of nucleic acids: a structure for deoxyribose nucleic acid'' \emph{Nature} \textbf{171} (1953) 737--738.
		
		\bibitem{8}
		M. Planck: ``Zur Theorie des Gesetzes der Energieverteilung im Normalspektrum'' \emph{Verhand. Deutschen Phys. Gesell.} 2 (1900) 237--245.
		
		\bibitem{9}
		N. Bohr: ``The quantum postulate and the recent development of atomic theory'' \emph{Nature} \textbf{121} (1928) 580--590. Reprinted in J. A. Wheeler and W. H. Zurek: \emph{Quantum Theory and Measurement} (Princeton University Press, 1983) p.87--126, referred to hereafter as WZ.
		
		\bibitem{10}
		N. Bohr: ``Can quantum-mechanical description of physical reality be considered complete?'' \emph{Phys. Rev.} \textbf{48} (1935) 696--702 reprinted in WZ, p. 145--151.
		
		\bibitem{11}
		A. Einstein to J.J. Laub, 1908, undated, Einstein Archives; scheduled for publication in \emph{The Collected Papers of Albert Einstein}, group of volumes on the Swiss years 1902--1914, Volume 5: \emph{Correspondence, 1902--1914} (Princeton University Press, Princeton, New Jersey).
		
		\bibitem{12}
		J.A. Wheeler: ``Assessment of Everett's `relative state' formulation of quantum theory'' \emph{Rev. Mod. Phys.} \textbf{29} (1957) 463--465.
		
		\bibitem{13}
		J.A. Wheeler: ``On the nature of quantum geometrodynamics,'' \emph{Ann. of Phys.} \textbf{2} (1957) 604--614.
		
		\bibitem{14}
		J. A. Wheeler: ``Superspace and the nature of quantum geometrodynamics,'' in \emph{Battelle Rencontres: 1967 Lectures in Mathematics and Physics}, eds. C. M. DeWitt and J. A. Wheeler, (Benjamin, New York, 1968) p. 242--307; reprinted as ``Le superspace et la nature de la g\'eometrodynamique quantique,'' in \emph{Fluides et Champ Gravitationnel en Relativit\'e G\'en\'erale}, No. 170, Colloques Internationaux (Editions de Centre National de la recherche Scientifique, Paris, 1969) p.257--322.
		
		\bibitem{15}
		J. A. Wheeler: ``Transcending the law of conservation of leptons,'' in \emph{Atti del Convegno Internazionale sul Tema: The Astrophysical Aspects of the Weak Interactions} (Cortona ``Il Palazzone,'' 10--12 Giugno 1970), Accademia Nazionale dei Lincei, Quaderno N. 157 (1971) p.133--64.
		
		\bibitem{16}
		C. W. Misner, K. S. Thorne and J. A. Wheeler: \emph{Gravitation} (Freeman, San Francisco, now New York, 1973) p. 1217, cited hereafter as MTW; paragraph on participatory concept of the universe.
		
		\bibitem{17}
		J. A. Wheeler: ``The universe as home for man,'' in \emph{The Nature of Scientific Discovery}, ed. O. Gingerich (Smithsonian Institution Press, Washington, 1975) p. 261--296; preprinted in part in \emph{American Scientist}, \textbf{62} (1974) 683--91; reprinted in part as T. P. Snow, \emph{The Dynamic Universe} (West, St. Paul, Minnesota, 1983) p.108--109.
		
		\bibitem{18}
		C. M. Patton and J. A. Wheeler: ``Is physics legislated by cosmogony?,'' in \emph{Quantum Gravity}, eds. C. Isham, R. Penrose and D. Sciama (Clarendon, Oxford, 1975) p. 538--605; reprinted in part in \emph{Encyclopaedia of Ignorance}, eds. R. Duncan and M. Weston-Smith (Pergamon, Oxford, 1977) p.19--35.
		
		\bibitem{19}
		J. A. Wheeler: ``Include the observer in the wave function?'' \emph{Fundamenta Scientiae: S\'eminaire sur les fondements des sciences} (Strasbourg) \textbf{25} (1976) 9--35; reprinted in \emph{Quantum Mechanics A Half Century Later}, eds. J. Leite Lopes and M. Paty (Reidel, Dordrecht, 1977) p.1--18.
		
		\bibitem{20}
		J. A. Wheeler: ``Genesis and observership,'' in \emph{Foundational Problems in the Special Sciences}, eds. R. Butts and J. Hintikka (Reidel, Dordrecht, 1977) p.1--33.
		
		\bibitem{21}
		J. A. Wheeler: ``The `past' and the `delayed choice' double-slit experiment,'' in \emph{Mathematical Foundations of Quantum Theory}, ed. A. R. Marlow (Academic, New York, 1978) p.9--48; reprinted in part in WZ, p.182--200.
		
		\bibitem{22}
		J. A. Wheeler: ``Frontiers of time,'' in \emph{Problems in the Foundations of Physics}, Proceedings of the International School of Physics ``Enrico Fermi'' (Course 72), ed. N. Toraldo di Francia (North-Holland, Amsterdam, 1979) p.395--497; reprinted in part in WZ, p. 200--208.
		
		\bibitem{23}
		J. A. Wheeler: ``The quantum and the universe,'' in \emph{Relativity, Quanta, and Cosmology in the Development of the Scientific Thought of Albert Einstein}, Vol. II, eds. M. Pantaleo and F. deFinis (Johnson Reprint Corp., New York, 1979) p.807--825.
		
		\bibitem{24}
		J. A. Wheeler: ``Beyond the black hole,'' in \emph{Some Strangeness in the Proportion: A Centennial Symposium to Celebrate the Achievements of Albert Einstein}, ed. H. Woolf (Addison-Wesley, Reading, Massachusetts, 1980) p.341--375; reprinted in part in WZ, p.208--210.
		
		\bibitem{25}
		J. A. Wheeler: ``Pregeometry: motivations and prospects,'' in \emph{Quantum Theory and Gravitation}, proceedings of a symposium held at Loyola University, New Orleans, May 23--26, 1979, ed. A. R. Marlow (Academic, New York, 1980) p.1--11.
		
		\bibitem{26}
		J. A. Wheeler: ``Law without law,'' in \emph{Structure in Science and Art}, eds. P. Medawar and J. Shelley (Elsevier North-Holland, New York and Excerpta Medica, Amsterdam, 1980) p.132--54.
		
		\bibitem{27}
		J. A. Wheeler: ``Delayed-choice experiments and the Bohr-Einstein dialog,'' in \emph{American Philosophical Society and the Royal Society: Papers read at a meeting, June 5, 1980} (American Philosophical Society, Philadelphia, 1980) p.9--40; reprinted in slightly abbreviated form and translated into German as ``Die Experimente der verzögerten Entscheidung und der Dialog zwischen Bohr und Einstein,'' in \emph{Moderne Naturphilosophie} ed. B. Kanitscheider (Königshausen and Neumann, Würzburg, 1984) p. 203--222; reprinted in \emph{Niels Bohr: A Profile}, eds. A. N. Mitra, L. S. Kothari, V. Singh and S. K. Trehan (Indian National Science Academy, New Delhi, 1985) p.139--168.
		
		\bibitem{28}
		J. A. Wheeler: ``Not consciousness but the distinction between the probe and the probed as central to the elemental quantum act of observation,'' in \emph{The Role of Consciousness in the Physical World}, ed. R. G. Jahn (Westview, Boulder, 1981) p. 87--111.
		
		\bibitem{29}
		J. A. Wheeler: ``The elementary quantum act as higgledy-piggledy building mechanism,'' in \emph{Quantum Theory and the Structures of Time and Space}, Papers presented at a Conference held in Tutzing, July, 1980, eds. L. Castell and C. F. von Weizsäcker (Carl Hanser, Munich, 1981) p.27--304.
		
		\bibitem{30}
		J. A. Wheeler: ``The computer and the universe,'' \emph{Int. J. Theo. Phys.} \textbf{21} (1982) 557--571.
		
		\bibitem{31}
		J. A. Wheeler: ``Bohr, Einstein, and the strange lesson of the quantum,'' in \emph{Mind in Nature}, Nobel Conference XVII, Gustavus Adolphus College, St. Peter, Minnesota, ed. Richard Q. Elvee (Harper and Row, New York, 1982) p.1--30 (also 88, 112, 113, 130--131, 148--40).
		
		\bibitem{32}
		J. A. Wheeler: \emph{Physics and Austerity} (in Chinese) (Anhui Science and Technology Publications, Anhui, China, 1982); reprinted in part (Lecture II), in \emph{Krisis}, Vol.1, No.2, ed. I. Masculescu (Klincksieck, Paris, 1983) p.671--75.
		
		\bibitem{33}
		J. A. Wheeler: ``Particles and geometry,'' in \emph{Unified Theories of Elementary Particles}, eds. P. Breitenlohner and H. P. Dürr (Springer, Berlin, 1982) p.189--217.
		
		\bibitem{34}
		J. A. Wheeler: ``Blackholes and new physics,'' in \emph{Discovery: Research and Scholarship at the University of Texas at Austin}, 7, No.2 (Winter 1982) 4--7.
		
		\bibitem{35}
		J. A. Wheeler: ``On recognizing law without law,'' \emph{Am. J. Phys.} \textbf{51} (1983) 398--404.
		
		\bibitem{36}
		J. A. Wheeler: ``Jenseits aller Zeitlichkeit,'' in \emph{Die Zeit}, Schriften der Carl Friedrich von Siemens-Stiftung, Vol. 6, eds. A. Peisl and A. Mohler (Oldenbourg, München, 1983) p.17--34.
		
		\bibitem{37}
		J. A. Wheeler: ``Elementary quantum phenomenon as building unit,'' in \emph{Quantum Optics, Experimental Gravitation, and Measurement Theory}, eds. P. Meystre and M. Scully (Plenum, New York and London, 1983) p.141--143.
		
		\bibitem{38}
		J. A. Wheeler: ``Bits, quanta, meaning,'' in \emph{Problems in Theoretical Physics}, eds. A. Giovannini, F. Mancini and M. Marinaro (University of Salerno Press, Salerno, 1984) p. 121--141; also in \emph{Theoretical Physics Meeting: Atti del Convegno, Amalfi, 6--7 maggio 1983} (Edizioni Scientifiche Italiane, Naples, 1984) p. 121--134; also in \emph{Festschrift in Honour of Eduardo R. Caianiello}, eds. A. Giovannini, F. Mancini, M. Marinaro, and A. Rimini (World Scientific, Singapore, 1989) p.133--154.
		
		\bibitem{39}
		J. A. Wheeler: ``Quantum gravity: the question of measurement,'' in \emph{Quantum Theory of Gravity}, ed. S. M. Christensen (Hilger, Bristol 1984) p.224--233.
		
		\bibitem{40}
		W. A. Miller and J. A. Wheeler: ``Delayed-choice experiments and Bohr's elementary quantum phenomenon,'' in \emph{Proceedings of International Symposium of Foundations of Quantum Mechanics in the Light of New Technology, Tokyo, 1983}, eds. S. Kamefuchi et al. (The Physical Society of Japan, Tokyo, 1984) p.140--151.
		
		\bibitem{41}
		J. A. Wheeler: ``Bohr's 'phenomenon' and 'law without law' '' in \emph{Chaotic Behavior in Quantum Systems}, ed. G. Casati (Plenum, New York, 1985) p.363--378.
		
		\bibitem{42}
		A. Kheyfets and J. A. Wheeler: ``Boundary of a boundary principle and geometric structure of field theories,'' \emph{Int. J. Theo. Phys.} \textbf{25} (1986) 573--580.
		
		\bibitem{43}
		J. A. Wheeler: ``Physics as meaning circuit: three problems,'' in \emph{Frontiers of Non-Equilibrium Statistical Physics}, eds. G. T. Moore and M. O. Scully (Plenum, New York, 1986) p.25--32.
		
		\bibitem{44}
		J. A. Wheeler: ``Interview on the role of the observer in quantum mechanics,'' in \emph{The Ghost in the Atom}, eds. P.C.W. Davies and J. R. Brown (Cambridge University Press, Cambridge, 1986) p.58--69.
		
		\bibitem{45}
		J. A. Wheeler: ``How come the quantum,'' in \emph{New Techniques and Ideas in Quantum Measurement Theory} ed. D. M. Greenberger (Ann. New York Acad. Sci. 480 (1987) p.304--316.)
		
		\bibitem{46}
		J. A. Wheeler: ``Hermann Weyl and the unity of knowledge,'' in \emph{Exact Sciences and their Philosophical Foundations}, eds. W. Deppert et al. (Lang, Frankfurt am Main, 1988) p.469--503; appeared in abbreviated form in \emph{American Scientist} \textbf{74} (1986) 366--375.
		
		\bibitem{47}
		J. A. Wheeler: ``World as system self-synthesized by quantum networking,'' \emph{IBM J. of Res. and Dev.} \textbf{32} (1988) 4--15; reprinted, in \emph{Probability in the Sciences}, ed. E. Agazzi (Kluwer, Amsterdam, 1988) p.103--129.
		
		\bibitem{48}
		D. M. Greenberger, ed.: \emph{New Techniques and Ideas in Quantum Measurement Theory} (Annals of the New York Academy of Sciences, Vol.480 (1986)).
		
		\bibitem{49}
		B. d'Espagnat: \emph{Reality and the Physicist: Knowledge, Duration and the Quantum World} (Cambridge University Press, Cambridge, 1989).
		
		\bibitem{50}
		P. Mittelstaedt and E. W. Stachow, eds: \emph{Recent Developments in Quantum Logic} (Bibliographisches Institut, Zürich, 1985).
		
		\bibitem{51}
		J. S. Bell: \emph{Speakable and Unspeakable in Quantum Mechanics: Collected Papers in Quantum Mechanics} (Cambridge University Press, Cambridge, 1987).
		
		\bibitem{52}
		J. W. Tukey: ``Sequential conversion of continuous data to digital data,'' Bell Laboratories memorandum of 1 September 1947 marks the introduction of the term ``bit'' reprinted in \emph{Origin of the term bit}, ed. H. S. Tropp (\emph{Annals Hist. Computing} 6 (1984)152--155.)
		
		\bibitem{53}
		W. K. Wootters and W. H. Zurek: ``A single quantum cannot be cloned,'' \emph{Nature} \textbf{279} (1982) 802--803.
		
		\bibitem{54}
		W. K. Wootters and W. H. Zurek, ``On replicating photons,'' \emph{Nature}, \textbf{304} (1983)188--189.
		
		\bibitem{55}
		Y. Aharonov and D. Bohm: ``Significance of electromagnetic potentials in the quantum theory'' \emph{Phys. Rev.} \textbf{115} (1959) 485--491.
		
		\bibitem{56}
		J. Anandan: ``Comment on geometric phase for classical field theories,'' \emph{Phys. Rev. Lett.} \textbf{60} (1988) 2555.
		
		\bibitem{57}
		J. Anandan and Y. Aharonov: ``Geometric quantum phase and angles,'' \emph{Phys. Rev.} D\textbf{38} (1988) 1863--1870.
		
		\bibitem{58}
		J. D. Bekenstein: ``Black holes and the second law,'' \emph{Nuovo Cimento Lett.} \textbf{4} (1972) 737--740.
		
		\bibitem{59}
		J. D. Bekenstein: ``Generalized second law of thermodynamics in black-hole physics,'' \emph{Phys. Rev.} D\textbf{9} (1973) 3292--3300.
		
		\bibitem{60}
		J. D. Bekenstein: ``Black-hole thermodynamics'' \emph{Physics Today} \textbf{33} (1980) 24--31.
		
		\bibitem{61}
		R. Penrose: ``Gravitational collapse: the role of general relativity,'' \emph{Riv. Nuovo Cimento} \textbf{1} (1969) 252--276.
		
		\bibitem{62}
		D. Christodoulou: ``Reversible and irreversible transformations in black-hole physics,'' \emph{Phys. Rev. Lett.} \textbf{25} (1970) 1596--1597.
		
		\bibitem{63}
		D. Christodoulou and R. Ruffini: ``Reversible transformations of a charged black hole,'' \emph{Phys. Rev.} D\textbf{4} (1971) 3552--3555.
		
		\bibitem{64}
		S.W. Hawking: ``Particle creation by black holes,'' \emph{Commun. Math. Phys.} \textbf{43} (1975) 199--220.
		
		\bibitem{65}
		S. W. Hawking: ``Black holes and thermodynamics,'' \emph{Phys. Rev.} \textbf{13} (1976) 191--197.
		
		\bibitem{66}
		W. H. Zurek and K. S. Thorne: ``Statistical mechanical origin of the entropy of a rotating, charged black hole,'' \emph{Phys. Rev. Lett.} \textbf{20} (1985) 2171--2175.
		
		\bibitem{67}
		MTW, p.1217.
		
		\bibitem{68}
		W. Shakespeare: \emph{The Tempest}, Act IV, Scene 1, lines 148 ff.
		
		\bibitem{69}
		T. Mann: \emph{Freud, Goethe, Wagner} (New York, 1937) p.20; trans. by H. T. Lowe-Porter from \emph{Freud und die Zukunft} (Vienna, 1936).
		
		\bibitem{70}
		G. W. Leibniz as cited in J. R. Newman: \emph{The World of Mathematics} (Simon and Schuster, New York, 1956).
		
		\bibitem{71}
		N. E. Steenrod: \emph{Cohomology Operations} (Princeton University Press, Princeton, New Jersey, 1962).
		
		\bibitem{72}
		C. Ehresmann: \emph{Catégories et Structures} (Dunod, Paris, 1965).
		
		\bibitem{73}
		D. Lohmer: \emph{Phänomenologie der Mathematik: Elemente einer Phänomenologischen Aufklärung der Mathematischen Erkenntnis nach Husserl} (Kluwer, Norwell, Massachusetts, 1989).
		
		\bibitem{74}
		A. Weil: ``De la métaphysique aux mathématiques,'' \emph{Sciences}, p.52--56; reprinted in A. Weil: \emph{Oeuvres Scientifiques: Collected Works}, Vol. 2, 1951--64 (Springer, New York, 1979) p.408--412.
		
		\bibitem{75}
		J. D. Barrow and F. J. Tipler: \emph{The Anthropic Cosmological Principle} (Oxford University Press, New York, 1986) and literature therein cited.
		
		\bibitem{76}
		See for example the survey by F. Feferman: ``Turing in the land of O(z),'' and related papers on mathematical logic, in R. Herken: \emph{The Universal Turing Machine: A Half Century Survey}, (Kammerer and Unverzagt, Hamburg and Oxford University Press, New York, 1988) p.113--147.
		
		\bibitem{77}
		H. Weyl: ``Mathematics and logic.'' A brief survey serving as a preface to a review of \emph{The Philosophy of Bertrand Russell}, \emph{Am. Math. Monthly} \textbf{53} (1946) 2--13.
		
		\bibitem{78}
		W. V. O. Quine: p.18 in the essay ``On what there is,'' in \emph{From a Logical Point of View}, 2nd ed. (Harvard University Press, Cambridge, Massachusetts, 1980) p.1--19.
		
		\bibitem{79}
		Discovered among graffiti in the men's room of the Pecan Street Cafe, Austin, Texas.
		
		\bibitem{80}
		G. W. Leibniz: \emph{Animadversiones ad Joh. George Wachteri librum de recondita Hebraeorum philosophia}, c. 1708, unpublished; English translation in P. P. Wiener, \emph{Leibniz Selections} (Scribners, New York, 1951) p.488.
		
		\bibitem{81}
		A. Einstein: as quoted by A. Forsee in \emph{Albert Einstein Theoretical Physicist} (Macmillan, New York, 1963) p.81.
		
		\bibitem{82}
		MTW, 43.4.
		
		\bibitem{83}
		Ref. 22, p.411.
		
		\bibitem{84}
		E. H. Spanier: \emph{Algebraic Topology} (McGraw-Hill, New York, 1966).
		
		\bibitem{85}
		E. Cartan: \emph{La Géométrie des Espaces de Riemann}, Memorial des Sciences Mathématiques (Gauthier-Villars, Paris, 1925).
		
		\bibitem{86}
		E. Cartan: \emph{Leçons sur la Géométrie des Espaces de Riemann} (Gauthier-Villars, Paris, 1925).
		
		\bibitem{87}
		MTW, Chap. 15.
		
		\bibitem{88}
		M. Atiyah: \emph{Collected Papers}. Vol. 5: \emph{Gauge Theories} (Clarendon, Oxford, 1988).
		
		\bibitem{89}
		Ref. 21, p.41--42; ref. 22, p.397--398.
		
		\bibitem{90}
		W. K. Wootters and W. H. Zurek: ``Complementarity in the double-slit experiment: quantum nonseparability and a quantitative statement of Bohr's principle,'' \emph{Phys. Rev.} D\textbf{19} (1979) 473--484.
		
		\bibitem{91}
		W. Heisenberg: ``Über den anschaulichen Inhalt der quantentheoretischen Kinematik und Mechanik,'' \emph{Zeits. f. Physik} \textbf{43} (1927) 172--198; English translation in WZ, p.62--84.
		
		\bibitem{92}
		N. Bohr and L. Rosenfeld: ``Zur Frage der Messbarkeit der elektromagnetischen Feldgrössen'' \emph{Mat.-fys. Medd. Dan. Vid. Selsk.} 12, no.8 (1933); English translation by Aage Petersen, 1979, reprinted in WZ, 479--534.
		
		\bibitem{93}
		D. J. Gross: ``On the calculation of the fine-structure constant,'' \emph{Phys. Today} 42, No.12 (1989).
		
		\bibitem{94}
		N. F. Mott: ``The wave mechanics of a-ray tracks,'' \emph{Proc. Roy. Soc. London} A\textbf{126} (1929) 74--84; reprinted in WZ, p.129--134.
		
		\bibitem{95}
		H. D. Zeh: ``On the interpretation of measurement in quantum theory,'' \emph{Found. Phys.} \textbf{1} (1970) 69--76.
		
		\bibitem{96}
		H. D. Zeh: \emph{The Physical Basis of the Direction of Time} (Springer, Berlin, 1989).
		
		\bibitem{97}
		E. Joos and H. D. Zeh: ``The emergence of classical properties through interaction with the environment,'' \emph{Zeits. f. Physik} B\textbf{59} (1985) 223--243.
		
		\bibitem{98}
		W. H. Zurek: ``Pointer basis of quantum apparatus: Into what mixture does the wavepacket collapse?,'' \emph{Phys. Rev.} D\textbf{24} (1981) 1516--1525.
		
		\bibitem{99}
		W. H. Zurek: ``Environment-induced superselection rules,'' \emph{Phys. Rev.} D\textbf{26} (1982) 1862--1880.
		
		\bibitem{100}
		W. H. Zurek: ``Information transfer in quantum measurements: irreversibility and amplification,'' in \emph{Quantum Optics, Experimental Gravitation and Measurement Theory}, eds. P. Meystre and M. O. Scully (Plenum, New York, 1983) p.87--116.
		
		\bibitem{101}
		W. G. Unruh and W. H. Zurek: ``Reduction of a wave packet in quantum Brownian motion,'' \emph{Phys. Rev.} D\textbf{40} (1989) 1071--1094.
		
		\bibitem{102}
		J. B. Hartle: ``Progress in quantum cosmology,'' preprint from Physics Department, University of California at Santa Barbara, 1989.
		
		\bibitem{103}
		M. B. Green, J. H. Schwarz and E. Witten: \emph{Superstring Theory} (Cambridge University Press, Cambridge, U.K., 1987).
		
		\bibitem{104}
		L. Brink and M. Henneaux: \emph{Principles of String Theory: Studies of the Centra de Estudios Cientificos de Santiago} (Plenum, New York, 1988).
		
		\bibitem{105}
		S. W. Hawking: ``The Boundary Conditions of the Universe,'' in \emph{Astrophysical Cosmology}, Pontificia Academic Scientiarum, eds. H. A. Brück, G. V. Coyne and M. S. Longair (Vatican City, 1982) p. 563--594.
		
		\bibitem{106}
		A. Vilenkin: ``Creation of universes from nothing,'' \emph{Phys. Lett.} B\textbf{117} (1982) 25--28.
		
		\bibitem{107}
		J. B. Hartle and S. W. Hawking: ``Wave function of the universe,'' \emph{Phys. Rev.} D\textbf{28} (1983) 2960--2975.
		
		\bibitem{108}
		W. K. Wootters: ``The acquisition of information from quantum measurements,'' Ph.D. dissertation, University of Texas at Austin (1980).
		
		\bibitem{109}
		W. K. Wootters: ``Statistical distribution and Hilbert space,'' \emph{Phys. Rev.} \textbf{23} (1981) 357--362.
		
		\bibitem{110}
		R. A. Fisher: ``On the dominance ratio,'' \emph{Proc. Roy. Soc. Edin.} \textbf{42} (1922) 321--341.
		
		\bibitem{111}
		R. A. Fisher: \emph{Statistical Methods and Statistical Inference} (Hafner, New York, 1956) p.8--17.
		
		\bibitem{112}
		B. C. G. Stueckelberg: ``Theoreme H et unitarité de S,'' \emph{Helv. Phys. Acta} \textbf{25} (1952) 577--580.
		
		\bibitem{113}
		E. C. G. Stueckelberg: ``Quantum theory in real Hilbert space,'' \emph{Helv. Phys. Acta} \textbf{33} (1960) 727--752.
		
		\bibitem{114}
		D. S. Saxon: \emph{Elementary Quantum Mechanics} (Holden, San Francisco 1964).
		
		\bibitem{115}
		H. J. Larson: \emph{Introduction to Probability Theory and Statistical Inference}, 2nd ed. (Wiley, New York, 1974).
		
		\bibitem{116}
		E. Schrödinger: ``The Foundation of the Theory of Probability,'' \emph{Proc. Roy. Irish Acad.} \textbf{51A} (1947) 51--66 and 141--146.
		
		\bibitem{117}
		E. T. Jaynes: ``Bayesian methods: General background,'' in \emph{Maximum Entropy and Bayesian Methods in Applied Statistics}, ed. J. H. Justice (Cambridge University Press, Cambridge, U.K., 1986) p.1--25.
		
		\bibitem{118}
		R. Viertl, ed.: \emph{Probability and Bayesian Statistics} (World Scientific, Singapore, 1987).
		
		\bibitem{119}
		R. D. Rosenkrantz, ed.: \emph{B. T. Jaynes: Papers on Probability, Statistics and Statistical Physics} (Reidel-Kluwer, Hingham, Massachusetts, 1989).
		
		\bibitem{120}
		P. J. Denning: ``Bayesian learning,'' \emph{American Sci.} \textbf{77} (1989) 216--218.
		
		\bibitem{121}
		J. O. Berger and D. A. Berry: ``Statistical analysis and the illusion of objectivity,'' \emph{American Sci.} \textbf{76} (1988) 159--165.
		
		\bibitem{122}
		J. Burke: \emph{The Day the Universe Changed} (Little, Brown, Boston, Massachusetts, 1985).
		
		\bibitem{123}
		F. Beck [pseudonym of the early nuclear-reaction-rate theorist Fritz Houtermans] and W. Godin: translated from the German original by E. Mosbacher and D. Porter, \emph{Russian Purge and The Extraction of Confessions} (Hurst and Blackett, London, 1951).
		
		\bibitem{124}
		D. Føllesdal: ``Meaning and experience,'' in \emph{Mind and Language}, ed. S. Guttenplan (Clarendon, Oxford, 1975) p. 25--44.
		
		\bibitem{125}
		G. Berkeley: \emph{Treatise Concerning the Principles of Understanding}, Dublin (1710; 2nd ed. 1734).
		
		\bibitem{126}
		T. Segerstedt: as quoted in ref. 22, p.415.
		
		\bibitem{127}
		Ref. 21, p.41.
		
		\bibitem{128}
		Ya. B. Zel'dovich and I. D. Novikov: \emph{Relativistic Astrophysics}, Vol. 1: \emph{Stars and Relativity} (University of Chicago Press, Chicago, 1971).
		
		\bibitem{129}
		J. Mather et al: ``A preliminary measurement of the cosmic microwave background spectrum by the Cosmic Background Explorer (COBE) Satellite,'' submitted for publication, \emph{Astrophys. J. Lett.} (1990).
		
		\bibitem{130}
		MTW, p.738, Box 27.4; or JGST, Chap. 13, p.242.
		
		\bibitem{131}
		G. K. O'Neill: \emph{The High Frontier}, 4th ed. (Space Studies Institute, Princeton, New Jersey, 1989).
		
		\bibitem{132}
		R. Jastrow: \emph{Journey to the Stars: Space Exploration --- Tomorrow and Beyond} (Bantam, New York, 1989).
		
		\bibitem{133}
		K. Popper: \emph{Conjectures and Refutations: the Growth of Scientific Knowledge} (Basic Books, New York, 1962).
		
		\bibitem{134}
		R. W. Fuller and P. Putnam: ``On the origin of order in behavior,'' \emph{General Systems} (Ann Arbor, Michigan) \textbf{12} (1966)111--121.
		
		\bibitem{135}
		R. W. Fuller: ``Causal and Moral Law: Their Relationship as Examined in Terms of a Model of the Brain,'' Monday Evening Papers (Wesleyan University Press, Middletown, Connecticut, 1967).
		
		\bibitem{136}
		G.M. Edelman: \emph{Neural Darwinism} (Basic Books, New York, 1987).
		
		\bibitem{137}
		W. H. Calvin: \emph{The Cerebral Symphony} (Bantam, New York, 1990).
		
		\bibitem{138}
		W. W. Collins: \emph{The Moonstone} (London, 1968).
		
		\bibitem{139}
		J. Allan Hobson: \emph{Sleep} (Scientific American Library, Freeman, New York, 1989) p.86, 89, 175, 185, 186.
		
		\bibitem{140}
		I. Langmuir: ``Pathological Science,'' 1953 colloquium, transcribed and edited, \emph{Phys. Today}, 42, No.12 (1989) 36--48.
		
		\bibitem{141}
		N. S. Hetherington: \emph{Science and Objectivity: Episodes in the History of Astronomy} (Iowa State University Press, Ames, Iowa, 1988).
		
		\bibitem{142}
		W. Sheehan: \emph{Planets and Perception: Telescopic Views and Interpretations} (University of Arizona Press, Tucson, Arizona, 1988).
		
		\bibitem{143}
		M. White: \emph{Science and Sentiment in America: Philosophical Thought from Jonathan Edwards to John Dewey} (Oxford University Press, New York, 1972).
		
		\bibitem{144}
		C. S. Peirce: \emph{The Philosophy of Peirce: Selected Writings}, ed. J. Buchler (Routledge and Kegan Paul, London, 1940), passages from p.337, 335, 336, 353 and 358; reprinted in ref. 18, p.593--595.
		
		\bibitem{145}
		Parmenides of Elea [c. 515 B.C.--450 B.C.], poem \emph{Nature}, part Truth, as summarized by A. C. Lloyd in article Parmenides, \emph{Encyclopedia Brittanica}, Chicago 17 (1959) 327.
		
		\bibitem{146}
		G.B. Pugh: \emph{On the Origin of Human Values} (New York, 1976); chapter Human values, free will, and the conscious mind preprinted in \emph{Zygon} \textbf{11} (1976) 2--24.
		
		\bibitem{147}
		F. W. J. von Schelling [1775--1854]: in \emph{Schellings Werke}, nach der Originalausgabe in neuer Anordnung herausgegeben, 6 vols., ed. M. Schröter, (Beck, München, 1958--1959), esp. Vol.5, p.428--430.
		
		\bibitem{148}
		J. von Neumann and O. Morgenstern: \emph{Theory of Games and Economic Behavior} (Princeton University Press, Princeton, New Jersey, 1944).
		
		\bibitem{149}
		J. Wang: \emph{Theory of Games} (Oxford University Press, New York, 1988).
		
		\bibitem{150}
		J. R. Pierce: \emph{Symbols, Signals and Noise: The Nature and Process of Communication} (Harper and Brothers, New York, 1961).
		
		\bibitem{151}
		M. Schwartz: \emph{Telecommunication Networks: Protocols, Modeling and Analysis} (Addison-Wesley, Reading, Massachusetts, 1987).
		
		\bibitem{152}
		M. S. Roden: \emph{Digital Communication Systems Design} (Prentice Hall, Englewood Cliffs, New Jersey, 1988).
		
		\bibitem{153}
		P. W. Anderson: ``More is different,'' \emph{Science} \textbf{177} (1972) 393--396.
		
		\bibitem{154}
		C. Mead and L. Conway: \emph{Introduction to VLSI Systems} (Addison-Wesley, Reading, Massachusetts, 1980).
		
		\bibitem{155}
		P. B. Schneck: \emph{Supercomputer Architecture} (Kluwer, Norwell, Massachusetts, 1987).
		
		\bibitem{156}
		F. E. Yates, ed.: \emph{Self-Organizing Systems: The Emergence of Order} (Plenum, New York, 1987).
		
		\bibitem{157}
		H. Haken: \emph{Information and Self-Organization: A Macroscopic Approach to Complex Systems} (Springer, Berlin, 1988).
		
		\bibitem{158}
		T. Kohonen: \emph{Self-Organization and Associative Memory}, 3rd ed. (Springer, New York, 1989).
		
		\bibitem{159}
		C. Smorynski: \emph{Self-reference and Model Logic} (Springer, Berlin, 1985).
		
		\bibitem{160}
		G. J. Chaitin: \emph{Algorithmic Information Theory}, rev. 1987 ed., (Cambridge University, Cambridge, 1988).
		
		\bibitem{161}
		J. P. Delahaye: ``Chaitin's equation; an extension of Gödel's theorem,'' \emph{Notices Amer. Math. Soc.} \textbf{36} (1989) 948--987.
		
		\bibitem{162}
		J. F. Traub, G. W. Wasilkowski and H. Wozniakowski: \emph{Information-Based Complexity} (Academic, San Diego, 1988).
		
		\bibitem{163}
		P. Young: \emph{The Nature of Information} (Praeger Greenwood, Westport, Connecticut, 1987).
		
		\bibitem{164}
		W. H. Zurek: ``Algorithmic randomness and physical entropy,'' \emph{Phys. Rev.} A\textbf{40} (1989) 4731--4751.
		
		\bibitem{165}
		M. Eigen and R. Winkler: \emph{Das Spiel: Naturgesetze steuern den Zufall} (Piper, München, 1975).
		
		\bibitem{166}
		W. M. Elsasser: \emph{Reflections on a Theory of Organisms} (Orbis, Frelighsburg, Quebec, 1987).
		
		\bibitem{167}
		G. Nicolis and I. Prigogine: \emph{Exploring Complexity: An Introduction} (Freeman, New York, 1989).
		
		\bibitem{168}
		S. Watanabe, ed.: \emph{Methodologies of Pattern Recognition} (Academic, New York, 1967).
		
		\bibitem{169}
		J. Tou and R. C. Gonzalez: \emph{Pattern Recognition Principles} (Addison-Wesley, Reading, Massachusetts, 1974).
		
		\bibitem{170}
		H. Haken, ed.: \emph{Pattern Formation by Dynamic Systems and Pattern Recognition} (Springer, Berlin, 1979).
		
		\bibitem{171}
		H. Small and E. Garfield: ``The geography of science: disciplinary and national mappings,'' \emph{J. of Info. Sci.} \textbf{11} (1985) 147--159.
		
		\bibitem{172}
		M. Agu: ``Field theory of pattern recognition,'' \emph{Phys. Rev.} A\textbf{37} (1988) 4415--4418.
		
		\bibitem{173}
		M. Minsky and S. Papert: \emph{Perceptrons: An Introduction to Computational Geometry}, 2nd ed. (Massachusetts Institute of Technology Press, Cambridge, Massachusetts, 1988).
		
		\bibitem{174}
		L. A. Steen: ``The science of patterns,'' \emph{Science} \textbf{240} (1988) 611--616.
		
		\bibitem{175}
		B. M. Bennett, D. D. Hoffman and C. Prakash: \emph{Observer Mechanics: A Formal Theory of Perception} (Academic, San Diego, California, 1989).
		
		\bibitem{176}
		J. A. Wheeler: ``Polyelectrons,'' \emph{Ann. New York Acad. Sci.} \textbf{46} (1946) 219--238.
		
		\bibitem{177}
		D. Bohm: ``The paradox of Einstein, Rosen and Podolsky,'' originally published as section 15-19, Chapter 22 of D. Bohm: \emph{Quantum Theory} (Prentice-Hall, Englewood Cliffs, N.J., 1950), reprinted in WZ, 356--368.
		
		\bibitem{178}
		W. H. Zurek: ``Thermodynamic cost of computation: Algorithmic complexity and the information metric,'' \emph{Nature} \textbf{34} (1989)119--124.
		
		\bibitem{179}
		E. F. Taylor and J. A. Wheeler: \emph{Spacetime Physics} (Freeman, San Francisco, 1963) p.102.
		
		\end{thebibliography}
	\end{document}

		
		
		

	
